% TMLR version: condensed Section 3.6. Full arXiv-v1 text in Appendix J.3.
In an exploratory analysis, the euclidean RDM on unit-centred features was decomposed into its norm
and angular components (Section~\ref{sec:gainmethods}; Fig.~\ref{fig:separation}); contrasts are
untrained minus trained at V1 and trained minus untrained at LOC. The identity held to a maximum
relative error of \SepOneMaxRelError{} (SEP1: met). SEP2 was not met: the V1 norm contrast
($F$ = center) was $+\SepOwnVoneNormCenter$ [$\SepOwnVoneNormCenterLo$, $\SepOwnVoneNormCenterHi$]
in the own CNN, with a CI including 0, and had a CI above 0 in \SepTwoModels{} of 6 ImageNet
models; under the pre-stated interpretation rule, the V1 ranking is therefore not attributed to the
norm or the angular component. On the unweighted angular component $A$, the trained networks were
ahead at V1 ($\SepOwnVoneAngleCenter$ [$\SepOwnVoneAngleCenterLo$, $\SepOwnVoneAngleCenterHi$] in
the own CNN; CI below 0 in \SepThreeModels{} of 6 ImageNet models; SEP3: met) and at LOC
($+\SepOwnLocAngleCenter$ [$\SepOwnLocAngleCenterLo$, $\SepOwnLocAngleCenterHi$]; CI above 0 in
\SepFourModels{} of 6; SEP4: met). The norm-weighted angular term $W$ ($+\WAngOwnVone$
[$\WAngOwnVoneLo$, $\WAngOwnVoneHi$] in the own CNN; interval excluding 0 in \WAngImnPosModels{}
of 6 ImageNet models) pointed in the direction of the norm component and was not resolved. Each
component accounted for about 0.1\% of the rank variance of the fMRI RDM (own CNN, V1, NATIVE
224\,px). Descriptively (no decision rule), the V1 angle contrast of the ImageNet models grew from
224 to 448\,px with a CI above 0 in \SepSlopeAngleCiPositive{} of 5 models, no norm slope had a CI
excluding 0, and in the own CNN neither component had a slope from 32 to 224\,px; convnext\_tiny
was an exception, with CIs including 0 for its V1 angle contrast and its angle slope
(Appendix~\ref{app:extgeometry}).
